% Part: proof-theory
% Chapter: cut-elimination
% Section: midsequent

\documentclass[../../../include/open-logic-section]{subfiles}

\begin{document}

\olfileid[id]{pt}{cut}{mid}

\olsection{Teorema Sekuen Tengah dan Teorema Herbrand}

Salah satu konsekuensi penting teorema eliminasi potongan adalah teorema
sekuen tengah. Teorema ini menyatakan bahwa jika terdapat
!!a{proof}~$\pi$ atas~$\Gamma \Sequent \Delta$ yang setiap
!!{formula}-nya berbentuk preneks, maka terdapat !!a{proof}~$\pi'$ yang
memuat suatu sekuen $S\ident\Pi\Sequent\Lambda$ (disebut
\emph{sekuen tengah}) sedemikian sehingga semua inferensi di atas $S$
dalam~$\pi'$ bersifat proposisional dan semua inferensi di bawah~$S$
merupakan inferensi kuantor. (!!^a{formula}~$!A$ disebut
\emph{preneks} (atau berada dalam \emph{bentuk preneks}) jika berbentuk
\[\mathsf{Q}_1x_1\dots\mathsf{Q}_nx_n\,!B(x_1,\dots,x_n)\]
dengan setiap $\mathsf{Q}_i$ berupa $\lforall$ atau~$\lexists$.)
Salah satu konsekuensi penting teorema sekuen tengah adalah teorema
Herbrand. Kasus umumnya agak sulit diuraikan, tetapi secara sangat umum
bentuknya adalah sebagai berikut: untuk setiap !!{sentence}~$!A$
logika orde pertama yang !!{provable}, terdapat tautologi
proposisional~$!A'$ yang darinya $!A$ dapat di-!!{prove}. Berikut versi
untuk !!{sentence} berbentuk~$\lexists[x][!B(x)]$, dengan $!B$ bebas
kuantor:

\begin{thm}[Teorema Herbrand]
Andaikan $!B(x)$ bebas kuantor dan
$\Proves\lexists[x][!B(x)]$. Maka terdapat suku-suku~$t_1$,
\dots, $t_n$ sedemikian sehingga
$\Proves !B(t_1)\lor\dots\lor !B(t_n)$.
\end{thm}

!!{formula}~$!B(t_1)\lor\dots\lor !B(t_n)$ disebut
\emph{disjungsi Herbrand} dari~$\lexists[x][!B(x)]$. Karena $!B$
bebas kuantor, disjungsi ini juga bebas kuantor. Jadi, jika disjungsi
tersebut !!{provable}, ia !!{provable} melalui bukti bebas potongan,
dan akibatnya dengan !!a{proof} yang hanya menggunakan inferensi
proposisional. Oleh karena itu, disjungsi tersebut merupakan tautologi.

Bagian sulit dari Teorema Herbrand ialah menunjukkan bahwa suatu disjungsi
Herbrand ada bagi setiap !!{formula}~$\lexists[x][!B(x)]$ yang
!!{provable}. Arah sebaliknya, yakni bahwa jika
$\lexists[x][!B(x)]$ memiliki disjungsi Herbrand maka formula tersebut
dapat dibuktikan, adalah mudah. Bahkan, disjungsi Herbrand itu secara logis mengimplikasikan
$\lexists[x][!B(x)]$. Misalnya, berikut !!a{proof} dalam~$\Log{G3c}$
untuk kasus $n=2$:
\[
\Axiom$!B(t_1) \fCenter \lexists[x][!B(x)], !B(t_1), !B(t_2)$
\Axiom$!B(t_2) \fCenter \lexists[x][!B(x)], !B(t_1), !B(t_2)$
\RightLabel{\LeftR{\lor}}
\BinaryInf$!B(t_1) \lor !B(t_2) \fCenter \lexists[x][!B(x)], !B(t_1), !B(t_2)$
\RightLabel{\RightR{\lexists}}
\UnaryInf$!B(t_1) \lor !B(t_2) \fCenter \lexists[x][!B(x)], !B(t_2)$
\RightLabel{\RightR{\lexists}}
\UnaryInf$!B(t_1) \lor !B(t_2) \fCenter \lexists[x][!B(x)]$
\DisplayProof
\]

Untuk mengekstrak disjungsi Herbrand dari !!a{proof}~$\pi$ atas
$\Sequent\lexists[x][!B(x)]$, mula-mula tinjau kasus ketika semua
inferensi $\RightR{\lexists}$ dalam !!{proof} bebas potongan~$\pi$
berada di bagian akhir:
\[
\AxiomC{}
\RightLabel{$\pi_1$}
\Deduce$ \fCenter \lexists[x][!B(x)], !B(t_1), \dots, !B(t_n)$
\RightLabel{\RightR{\lexists}}
\UnaryInf$ \fCenter \lexists[x][!B(x)], !B(t_2), \dots, !B(t_n)$
\RightLabel{$(n-1) \times \RightR{\lexists}$}
\Deduce$ \fCenter \lexists[x][!B(x)]$
\DisplayProof
\]
Jika semua inferensi $\RightR{\lexists}$ dalam~$\pi$ dengan
$\lexists[x][!B(x)]$ aktif adalah inferensi-inferensi ini, maka tidak
ada inferensi kuantor lain dalam~$\pi_1$. Alasannya, $!B(x)$ tidak
memuat kuantor lain, dan sekuen akhir tidak memuat !!{formula} apa pun
di sebelah kiri. Dengan kata lain,
$\fCenter\lexists[x][!B(x)],!B(t_1),\dots,!B(t_n)$ merupakan sekuen
tengah dari~$\pi$. Selain itu, karena tidak ada inferensi kuantor di atas
sekuen tengah, semua sekuen di atasnya memuat
$\lexists[x][!B(x)]$, dan kemunculan-kemunculan ini tidak aktif maupun
utama. Maka kemunculan-kemunculan itu dapat dihapus dari
!!{proof}~$\pi_1$, dan kita memperoleh !!a{proof} atas
$\fCenter !B(t_1),\dots,!B(t_n)$. Dengan menggunakan $n-1$ penerapan
$\RightR{\lor}$, kita lalu memperoleh !!a{proof} atas disjungsi
Herbrand~$\Sequent !B(t_1)\lor\dots\lor !B(t_n)$.

Masalahnya, kita tidak dapat mengasumsikan bahwa bahkan suatu !!{proof}
bebas potongan atas $\Sequent\lexists[x][!B(x)]$ memiliki semua
inferensi \RightR{\lexists} dalam satu blok di bagian akhir. Di antara
inferensi-inferensi \RightR{\lexists} itu mungkin terdapat inferensi
yang menjadikan salah satu $!B(t_i)$ sebagai formula utama. Di sinilah
teorema sekuen tengah diperlukan.

\begin{thm}[Teorema Sekuen Tengah]\ollabel{thm:midsequent}
Jika terdapat \Log{G3c}-!!{proof}~$\pi$ atas
$\Gamma\Sequent\Delta$, dengan semua !!{formula} dalam $\Gamma$ dan
$\Delta$ berbentuk preneks, maka terdapat !!a{proof}~$\pi'$ yang memuat
suatu sekuen $\Pi\Sequent\Lambda$ sedemikian sehingga semua inferensi di
bawahnya merupakan inferensi kuantor dan semua inferensi di atasnya
merupakan inferensi proposisional.
\end{thm}

\begin{proof}
  Definisikan \emph{orde} suatu inferensi kuantor dalam~$\pi$ sebagai
  banyaknya inferensi proposisional di bawahnya, dan orde~$o(\pi)$ dari
  $\pi$ sebagai jumlah orde semua inferensi kuantornya. Jika
  $o(\pi)=0$, tidak ada inferensi kuantor yang memiliki inferensi
  proposisional di bawahnya, dan kita selesai: karena semua inferensi
  kuantor mempunyai satu premis, terdapat tepat satu inferensi kuantor
  paling atas. Premisnya adalah sekuen tengah.

  Jika $o(\pi)>0$, pilih inferensi kuantor paling bawah yang berorde
  $>0$. Karena ordenya $>0$, di bawahnya pasti terdapat inferensi
  proposisional. Karena inferensi kuantor itu merupakan inferensi
  semacam ini yang paling bawah, kesimpulannya tidak mungkin menjadi
  premis suatu inferensi kuantor: jika demikian, inferensi kuantor kedua
  yang lebih bawah itu juga akan memiliki inferensi proposisional di
  bawahnya. Kita membedakan (banyak sekali!) kasus menurut inferensi
  kuantor yang dihadapi dan inferensi proposisional di bawahnya.
  Karena sekuen akhir hanya terdiri atas formula preneks, formula utama
  inferensi kuantor tidak dapat aktif dalam inferensi proposisional
  berikutnya. Jika aktif, formula utama inferensi proposisional itu akan
  memiliki kuantor dalam lingkup suatu penghubung proposisional.
  Berdasarkan sifat sub-!!{formula}, formula tersebut harus merupakan
  sub-!!{formula} dari !!a{formula} dalam sekuen akhir. Jadi, !!{formula}
  utama inferensi kuantor merupakan !!a{element} dari konteks inferensi
  proposisional yang lebih bawah.

  Berikut dua contoh.

Pertama, andaikan terdapat suatu $\RightR{\lforall}$ yang berakhir pada
premis kanan $\LeftR{\lif}$. Maka sub-!!{proof} terkait dari~$\pi$
berakhir pada sub-!!{proof}~$\pi_1$:
\[
\AxiomC{}
\RightLabel{$\pi_2$}
\Deduce$\Gamma' \fCenter \Delta', \lforall[x][!A(x)], !B$
\AxiomC{}
\RightLabel{$\pi_3$}
\Deduce$!C, \Gamma' \fCenter \Delta', !A(c)$
\RightLabel{\RightR{\lforall}}
\UnaryInf$!C, \Gamma' \fCenter \Delta', \lforall[x][!A(x)]$
\RightLabel{\LeftR{\lif}}
\BinaryInf$!B \lif !C, \Gamma' \fCenter \Delta', \lforall[x][!A(x)]$
\DisplayProof
\]
Kita dapat mengasumsikan bahwa $\pi$ reguler, sehingga variabel
eigen~$c$ tidak muncul di luar~$\pi_3$; khususnya, $c$ tidak muncul
dalam $!B\lif !C$ atau~$\pi_2$. Sekarang tinjau bukti~$\pi_1'$:
\[
\AxiomC{}
\RightLabel{$\pi_2'$}
\Deduce$\Gamma' \fCenter \Delta', \lforall[x][!A(x)], !A(c), !B$
\AxiomC{}
\RightLabel{$\pi_3'$}
\Deduce$!C, \Gamma' \fCenter \Delta', \lforall[x][!A(x)], !A(c)$
\RightLabel{\LeftR{\lif}}
\BinaryInf$!B \lif !C, \Gamma' \fCenter \Delta', \lforall[x][!A(x)], !A(c)$
\RightLabel{\RightR{\lforall}}
\UnaryInf$!B \lif !C, \Gamma' \fCenter \Delta', \lforall[x][!A(x)], \lforall[x][!A(x)]$
\DisplayProof
\]
Di sini $\pi_2'$ diperoleh dari~$\pi_2$ melalui pelemahan dengan~$!A(c)$
di sebelah kanan, dan $\pi_3'$ diperoleh dari~$\pi_3$ melalui pelemahan
dengan~$\lforall[x][!A(x)]$ di sebelah kanan. Pelemahan ini tidak
menambahkan inferensi apa pun; khususnya, tidak ada inferensi kuantor
yang dapat memengaruhi orde. Pelemahan tersebut juga tidak melanggar
syarat variabel eigen mana pun dalam~$\pi_2$ atau~$\pi_3$. Syarat
variabel eigen bagi \RightR{\lforall} tetap terpenuhi karena $c$ tidak
muncul dalam~$!B\lif !C$.

Kita menggunakan dapat diterimanya kontraksi untuk memperoleh
!!a{proof}~$\pi_1''$ atas
$!B\lif !C,\Gamma'\fCenter\Delta',\lforall[x][!A(x)]$. Bukti dapat
diterimanya $\RightR{\Contraction}$ bagi
$\lforall[x][!A(x)]$ dan bukti keterbalikan $\RightR{\lforall}$ yang
digunakan di dalamnya tidak mengubah orde inferensi apa pun: paling
jauh, bukti-bukti itu menghapus inferensi. Jadi, jumlah inferensi kuantor
dalam~$\pi_1''$ tidak lebih banyak daripada dalam~$\pi_1$, dan setiap
inferensi kuantor dalam~$\pi_1''$ memiliki jumlah inferensi proposisional
di bawahnya yang sama dengan inferensi pasangannya dalam~$\pi_1$---kecuali
\RightR{\lforall} terakhir. Dalam~$\pi$, inferensi itu memiliki
\LeftR{\lif} di bawahnya, sedangkan dalam~$\pi_1''$ kita telah
memindahkan \LeftR{\lif} ke atasnya. Ganti $\pi_1$ dalam $\pi$ dengan
$\pi_1''$. !!{proof} yang dihasilkan memiliki orde lebih rendah, karena
setidaknya orde inferensi $\RightR{\lforall}$ telah berkurang
sekurang-kurangnya~$1$. Sekarang terapkan hipotesis induksi.

Kasus ketika inferensi kuantornya adalah \RightR{\lexists} bahkan lebih
mudah, karena kontraksi tidak diperlukan dan tidak ada syarat variabel
eigen yang harus diperhatikan. Misalnya, andaikan $\RightR{\lexists}$
diikuti oleh $\RightR{\land}$ (kali ini sebagai premis kiri).
Sub-!!{proof} terkait adalah $\pi_1$:
\[
\AxiomC{}
\RightLabel{$\pi_2$}
\Deduce$\Gamma' \fCenter \Delta', \lexists[x][!A(x)], !A(t), !B$
\RightLabel{\RightR{\lexists}}
\UnaryInf$\Gamma' \fCenter \Delta', \lexists[x][!A(x)], !B$
\AxiomC{}
\RightLabel{$\pi_3$}
\Deduce$\Gamma' \fCenter \Delta', \lexists[x][!A(x)], !C$
\RightLabel{\RightR{\land}}
\BinaryInf$\Gamma' \fCenter \Delta', \lexists[x][!A(x)], !B \land !C$
\DisplayProof
\]
Ganti $\pi_1$ dalam~$\pi$ dengan $\pi_1'$:
\[
\AxiomC{}
\RightLabel{$\pi_2$}
\Deduce$\Gamma' \fCenter \Delta', \lexists[x][!A(x)], !A(t), !B$
\AxiomC{}
\RightLabel{$\pi_3'$}
\Deduce$\Gamma' \fCenter \Delta', \lexists[x][!A(x)], !A(t), !C$
\RightLabel{\RightR{\land}}
\BinaryInf$\Gamma' \fCenter \Delta', \lexists[x][!A(x)], !A(t), !B \land !C$
\RightLabel{\RightR{\lexists}}
\UnaryInf$\Gamma' \fCenter \Delta', \lexists[x][!A(x)], !B \land !C$
\DisplayProof
\]
Di sini $\pi_3'$ diperoleh dari $\pi_3$ melalui pelemahan dengan~$!A(t)$
di sebelah kanan. Pelemahan ini hanya menambahkan $!A(t)$ ke sebelah
kanan setiap sekuen dalam~$\pi_3$, sehingga tidak mengubah orde inferensi
kuantor apa pun. Orde inferensi kuantor dalam~$\pi_1'$ sama dengan
ordenya dalam~$\pi_1$, kecuali orde inferensi \RightR{\lexists} yang
dipertukarkan dengan~\RightR{\land} telah berkurang~$1$. Hipotesis
induksi berlaku.
\end{proof}

\begin{prob}
Uraikan kasus-kasus dalam bukti \olref{thm:midsequent} ketika
\LeftR{\lexists} berakhir pada premis~\RightR{\lif}, dan ketika
\LeftR{\lforall} berakhir pada premis kiri~\LeftR{\lor}.
\end{prob}

\begin{proof}[Bukti Teorema Herbrand]
Andaikan $\pi$ berakhir pada $\Sequent\lexists[x][!A(x)]$.
Berdasarkan \olref{thm:midsequent}, kita dapat mengubah $\pi$ menjadi
!!a{proof}~$\pi'$ dengan sekuen tengah~$S$, yang harus berbentuk
\[\Sequent\lexists[x][!A(x)],!A(t_1),\dots,!A(t_n).\]
Karena semua inferensi di atas $S$ bersifat proposisional,
$\lexists[x][!A(x)]$ tidak dapat menjadi formula utama inferensi di
atas~$S$. Karena formula itu tidak atomik, ia tidak dapat menjadi
formula utama dalam aksioma. Karena $\lexists[x][!A(x)]$ tidak dapat
menjadi subformula wajar dari $!A(t_1)$ (ingat bahwa $!A(x)$ bebas
kuantor), formula itu juga tidak dapat aktif di atas~$S$. Oleh sebab
itu, menghapus $\lexists[x][!A(x)]$ dari sub-!!{proof} yang berakhir
pada~$S$ menghasilkan !!{proof} yang benar atas
$\Sequent !A(t_1),\dots,!A(t_n)$. Dari bukti itu, !!a{proof} atas
$\Sequent !A(t_1)\lor\dots\lor !A(t_n)$ dapat diperoleh dengan
$n-1$ inferensi $\RightR{\lor}$.
\end{proof}

% \begin{prob}
% Find a Herbrand disjunction of $\lexists[x][\lforall[y][(!A(x) \lif
% !A(y))]]$ by finding a cut-free !!{proof} of $\Sequent
% \lexists[x][\lforall[y][(!A(x) \lif !A(y))]]$ and applying the
% transformation in \olref{thm:midsequent} to find a midsequent (if
% necessary).
% \end{prob}

\end{document}
